    \documentclass[a4paper,12pt]{article}

% *** FUENTE ***
\usepackage[T1]{fontenc}
\usepackage[lining,light]{InriaSans}
\renewcommand*\oldstylenums[1]{{\fontfamily{InriaSans-OsF}\selectfont #1}}
%% The font package uses mweights.sty which has som issues with the
%% \normalfont command. The following two lines fixes this issue.
\let\oldnormalfont\normalfont
\def\normalfont{\oldnormalfont\mdseries}

% *** COLOR ***
\usepackage{xcolor}
\definecolor{rojo}{RGB}{217, 62, 37}

\begin{document}
\begin{center}
    % Poner el título importante
    {\color{rojo} \LARGE \textbf {Titulo importante y poderoso}} \\[1cm]

\hspace{3cm } {\large \textit{Mi proyecto pro en LaTex}}

\underline{Chuck Norris}
\end{center}

\begin{flushright}
    19 de febrero del 2022
\end{flushright}

TeX1 \newline
 
\indent is a low-level markup
\noindent and programming language created by Donald Knuth2 to
typeset documents attractively and consistently. Knuth started writing the TeX typesetting
engine in 1977 to explore the potential of the digital printing equipment that was beginning
to infiltrate the publishing industry at that time, especially in the hope that he could reverse
the trend of deteriorating typographical quality that he saw affecting his own books and
articles. With the release of 8-bit character support in 1989, TeX development has been
essentially frozen with only bug fixes released periodically. TeX is a programming language
in the sense that it supports the if-else construct: you can make calculations with it (that
are performed while compiling the document), etc., but you would find it very hard to do
anything else but typesetting with it. The fine control TeX offers over document structure
and formatting makes it a powerful—and formidable—tool.

\textbackslash hola, 100 \% adas. \$ \# \& \{hola\} hola\_2
 
\end{document}